\scanpage{228}\mspage{83}
va faire au § suivant. Le cas que je vais y
traiter est le seul de son espèce que j'aie
\deleted{\unsure{co}}rencontré -- un cas de ``dualité formelle'' qui
s'introduit de façon naturelle, et qui
\unsure{ne} soit ``de $1^{\text{\unsure{ière|ère}}}$ espèce''.

6. \uline{Transporteurs}, \uline{noyaux}, \uline{stabiliseurs relatifs à} $W$.

Soit $M$ une catégorie, munie d'une partie
\[
W\subset\operatorname{Fl}(M).\tag*{(6.1)}
\]
Je pense au cas d'un localiseur, mais n'introduirai d'hypothèse sur $W$ qu'au besoin.
Pour simplifier, je suppose $M$ stable
par produits fibrés.

Je pose \margin{cf rectification p.~86}
\[
\left\{\begin{aligned}
P&=\unsure{\mathfrak{P}|\mathcal{P}}(W)=\{A\in\mathfrak{P}(\operatorname{Fl}(M))\mid A\subset W\}\\
Q&=\{L\in\mathfrak{P}(\operatorname{Fl}(M))\mid L\ \text{\uline{stable} par chgt de}\ \}\\[-2pt]
 &\qquad\text{base }\unsure{(}\text{par quoi je sous-entends}\\[-2pt]
 &\qquad\text{aussi, stable par}\\[-2pt]
 &\qquad\text{isomorphie) et par}\\[-2pt]
 &\qquad\text{composition}.
\end{aligned}\right.\tag*{(6.2)}
\]

J'introduis une relation
\[
\mathcal{R}\subset P\times Q
\]
de \unsure{fa|la|\deleted{fa}} façon suivante. Pour $A\in P$, $L\in Q$,
\unsure{donc|on} pose
\diagram{\begin{tikzcd}[ampersand replacement=\&,column sep=small,row sep=small]
{} \& {} \& {\text{Pour \unsure{tout|tt} diagramme cartésien}} \& {} \\
{\text{(6.3)}\quad(A,L)\in\mathcal R\ \Longleftrightarrow} \& X\arrow[d,"f"] \& X'\arrow[l]\arrow[d,"\unsure{f'|g'}"] \& {\text{où }f\in A,\ p\in L,}\\
{} \& Y \& Y'\arrow[l,"p"] \& {\text{on a } f'\in\deleted{\unsure{A}}W}
\end{tikzcd}}
(i.e. \deleted{\unsure{le chgt de base par un} $p\in L$}
\inserted{$A$ est stable par chgt de base}
\deleted{\unsure{d'un}} \inserted{par} $\unsure{p|f}\in\unsure{L|A}$)
\deleted{\unsure{donne un} $f'\in W$)}.
% archive sheet number: 228
